\begin{frame}{유명한 사례: 딥러닝 이전 시대, 우편번호를 읽던 SVM}
  \begin{itemize}
    \item 1990년대 미국 우편 서비스(USPS)는 \textbf{손으로 쓴
      우편번호}를 자동으로 읽어내는 문제로 골머리를 앓았다.
    \item 초기 신경망(LeCun의 CNN 포함)과 나란히 가장 강력한
      성능을 보인 방법 중 하나가 바로 \textbf{RBF 커널 SVM} —
      필기체 숫자처럼 명확한 직선으로 나뉘지 않는 복잡한 패턴에서도
      오차율을 \textbf{한 자릿수 퍼센트대}까지 끌어내림.
    \item 2010년대 초반 딥러닝(특히 10장 CNN, Krizhevsky et al.
      2012)이 대세가 되기 전까지 커널 SVM은 이미지 인식부터
      텍스트 분류(스팸 필터링에도 나이브베이즈와 나란히),
      생물정보학의 유전자 데이터 분류까지 폭넓게 쓰인
      ``\textbf{만능 분류기}''.
  \end{itemize}
  \vskip0.5em
  \begin{block}{지금도}
    데이터가 \textbf{수천$\sim$수만 개 수준}으로 많지 않고 특징 수가
    적당할 때는, 신경망보다 \textbf{학습이 훨씬 빠르고
    하이퍼파라미터도 적은} 커널 SVM이 실무에서 여전히 합리적인
    선택지.
  \end{block}
\end{frame}
